% !PRINTED 517
\weberpage{517}

%<<CONT
bezeichnen, die die einzelnen Reihen $A, B \dots S$ an ihrer
Stelle lassen und nur die Indices in jeder Reihe permutiren.
Diese Permutationen bilden für sich eine Gruppe, da durch die
Zusammensetzung von zweien unter ihnen die Reihen $A, B \dots
S$ nicht vertauscht werden können. Diese Gruppe $Q$ ist ein
\gesperrt{Normaltheiler} von $P$. Denn ist $\varkappa$ eine Permutation von
$Q$, und $\pi$ eine aus $P$, so gehört die zusammengesetzte
Permutation $\pi^{-1} \varkappa \pi$ wieder zu $Q$, weil die
Permutation unter den $A, B \dots S$, die durch $\pi^{-1}$
hervorgerufen ist, durch $\varkappa$ nicht geändert und durch
$\pi$ wieder rückgängig gemacht wird. Die Gruppe $Q$ ist
intransitiv, und durch Adjunction einer zu $Q$ gehörigen
Function $\psi$ wird nicht nur die \gesperrt{Galois}'sche Resolvente
reducibel, sondern die Function $f(x)$ selbst zerfällt in $s$
Factoren vom $r^{\mathrm{ten}}$ Grade,

\weq{2}{f(x) = f_\alpha(x,\psi)\, f_\beta(x,\psi) \dots f_\sigma(x,\psi).}

Die Wurzeln von $f_\alpha(x,\psi) = 0$ sind die Grössen $\alpha$,
und die \gesperrt{Galois}'sche Gruppe dieser Gleichung besteht aus allen
Permutationen, die durch $Q$ unter den $\alpha$ hervorgerufen
werden.

Die Grösse $\psi$ ist, wenn $j$ der Index des Theilers $Q$ von
$P$ ist, Wurzel einer Partialresolvente $j^{\mathrm{ten}}$
Grades $\chi(u) = 0$, die eine Normalgleichung und also ihre
eigene \gesperrt{Galois}'sche Resolvente ist (\S.\,155). Bezeichnen wir
ihre Wurzeln mit

\weq{3}{\psi,\ \psi_1,\ \psi_2 \dots \psi_{j-1},}

so besteht die \gesperrt{Galois}'sche Gruppe dieser Resolvente aus den
Substitutionen

\weq{4}{(\psi,\psi),\ (\psi,\psi_1),\ (\psi,\psi_2) \dots (\psi,\psi_{j-1}),}

die durch die Permutationen der Nebengruppen

\weq{5}{Q,\ Q\pi_1,\ Q\pi_2 \dots Q\pi_{j-1}}

hervorgerufen werden. Jede dieser Nebengruppen bringt eine
gewisse Permutation unter den Reihen $A, B \dots S$ hervor, und
zwar sind diese Permutationen alle von einander verschieden;
denn wenn etwa durch $\pi_1$ und $\pi_2$ dieselbe Permutation
unter den Reihen hervorgerufen wird, so gehört $\pi_2 \pi_1^{-1}$
zu $Q$ und $Q\pi_1$, $Q\pi_2$ sind nicht verschiedene
Nebengruppen. Die Permutationsgruppe der Partialresolvente
$\chi(u) = 0$ ist also mit der Gruppe der Permutationen, die
durch $P$ unter den Reihen $A, B \dots S$ hervorgerufen werden,
isomorph.

Verstehen wir unter $Q_\alpha$ den Inbegriff der Permutationen
von $P$, durch die die Reihe $A$ nicht verschoben wird, so ist
$Q_\alpha$
%CONT>>
